\documentclass[10pt,a4paper]{amsart}
\usepackage[margin=19mm]{geometry}
\usepackage{amsmath,amssymb,mathtools}
\usepackage[scr=rsfs,cal=euler]{mathalfa}
\usepackage{xfrac}
\usepackage[unicode,hidelinks,bookmarks=false]{hyperref}
\newcommand{\ie}{i.e.\ }
\newcommand{\cf}{cf.\ }
\newcommand{\resp}{respectively\ }
\newcommand{\Subsection}{\subsection}
\providecommand{\nobreakdash}{\nobreak\hbox{-}}
\setlength{\emergencystretch}{2em}
\multlinegap0pt
\usepackage{amssymb}
\newtheorem{theorem}{Theorem}
\newtheorem{proposition}{Proposition}
\title{A phase transition in the directional growth\\ of lamplighter groups}
\author{Mohammad F. Marashdeh}
\date{Source: arXiv:2607.21658v1 (CC BY 4.0)}
\begin{document}
\maketitle

\noindent\emph{Source and licence.} A phase transition in the directional growth\\ of lamplighter groups. Version arXiv:2607.21658v1 (CC BY 4.0), distributed by arXiv under the Creative Commons Attribution 4.0 International licence, http://creativecommons.org/licenses/by/4.0/, as recorded in the arXiv licence metadata for this exact version and shown on the abstract page https://arxiv.org/abs/2607.21658v1. Full licence notice: \texttt{licenses/arxiv-2607.21658-CC-BY-4.0.txt}.\par\smallskip

\noindent\textit{Editorial scope.} Counted unit: the article's proposition prop:variational (source0.tex lines 224--232). The article's complete original proof of it is delivered with it (source0.tex lines 234--277, one \texttt{proof} environment of 44 lines). No mathematics is written, repaired or re-derived: the statement and the proof are the article's own lines, in the article's own notation, and no retained line is substituted or re-flowed.  Retained uncounted as the source-local context the proof needs: lines 76--87.  Nothing was omitted from any retained span, so the package carries no omission record.  No retained line contains a citation, so the wrapper carries no bibliography and no bibliography entry.  No retained line contains a figure environment, an image-inclusion command or drawing code, so no image is delivered and the package declares no asset dependency.  Editorial formatting only, in addition: an AMS \texttt{amsart} preamble that restates the article's own declarations and macros used by the retained lines, namely the environments \texttt{theorem}, \texttt{proposition} on separate counters, together with the standard packages those lines need.  The retained environments print in this document as Theorem 1, Proposition 1 in order of appearance, whereas the article prints the same items with the numbering of its own full document.  The complete native source of the article as distributed in the arXiv e-print for this exact version is bundled unchanged as \texttt{sources/205885/source0.tex}.
\begin{theorem}[Spectrum]\label{thm:spectrum}
For every $\beta\in[-1,1]$ and every sequence of integers $j_n$ with $j_n/n\to\beta$, the limit
$I(\beta)=\lim_n\frac1n\log N_{j_n}(n)$ exists, is even and concave, and equals
\[
I(\beta)\;=\;
\begin{cases}
(1+|\beta|)\,\log\rho, & |\beta|\le\beta^{*},\\[2pt]
|\beta|\;h\!\Big(\dfrac{1-|\beta|}{|\beta|}\Big), & \beta^{*}\le|\beta|\le1 .
\end{cases}
\]
The same formula holds for the sphere-fiber counts $N_j(n)-N_j(n-1)$.
\end{theorem}

\begin{proposition}\label{prop:variational}
For every $\beta\in[0,1]$, $\ \Psi(\beta)=I(\beta)$ as defined in Theorem \ref{thm:spectrum}, the maximizer is unique, and it equals
\[
(\ell_\beta,\mu_\beta)=
\Big(\tfrac{1+\beta}{2+u_r},\ \tfrac{(1+\beta)u_r}{2+u_r}\Big)\ \ (\beta\le\beta^{*}),
\qquad
(\ell_\beta,\mu_\beta)=\big(\beta,\ 1-\beta\big)\ \ (\beta\ge\beta^{*}).
\]
\end{proposition}

\begin{proof}
\emph{Positivity.} Let $\beta\in[0,1)$ and pick any $u\in(0,u_{\max}(\beta)]$, where $u_{\max}(\beta)=\min\big(1,(1-\beta)/\beta\big)$ and $u_{\max}(0):=1$. Setting $\ell=(1+\beta)/(2+u)$ and $\mu=u\ell$ gives $\mu+2\ell-\beta=\ell(2+u)-\beta=1$ and $\ell\ge\beta$, so $(\ell,\mu)\in \Delta_\beta$ and $\Psi(\beta)\ge\ell h(u)>0$.

\emph{The budget binds.} For $0<u\le1$,
\begin{equation}\label{eq:partials}
\frac{\partial}{\partial\ell}\,\ell h(\mu/\ell)\Big|_{\mu}\;=\;h(u)-u h'(u)\;=\;-\log(1-u)\;>\;0,
\qquad u=\mu/\ell ,
\end{equation}
interpreted as $+\infty$ at $u=1$. Let $(\ell,\mu)$ maximize. If $\mu=0$ the value is $0$, which by positivity is impossible for $\beta<1$; and for $\beta=1$ the constraints $\ell\ge1$ and $\mu+2\ell\le2$ force $\ell=1$, $\mu=0$, where the budget binds. So assume $\mu>0$. If $\mu+2\ell-\beta<1$, increasing $\ell$ slightly keeps $\ell\ge\beta$ and $\mu\le\ell$, keeps the budget satisfied, and strictly increases the value by \eqref{eq:partials}, a contradiction. Hence $\mu+2\ell-\beta=1$ at every maximizer.

\emph{Reduction to one variable.} On the budget face, writing $u=\mu/\ell\in[0,1]$ gives $\ell(2+u)=1+\beta$, so the constraint $\ell\ge\beta$ reads $\beta(2+u)\le1+\beta$, that is $u\le(1-\beta)/\beta$, vacuous at $\beta=0$. Therefore
\[
\Psi(\beta)=(1+\beta)\,\max_{0\le u\le u_{\max}(\beta)}\psi(u),
\qquad
\psi(u)=\frac{h(u)}{2+u}.
\]

\emph{The unconstrained maximum of $\psi$.} Let $g(u)=(2+u)h'(u)-h(u)$, so $\psi'(u)=g(u)/(2+u)^2$. Then $g'(u)=(2+u)h''(u)<0$, while $h'(u)=\log\frac{r(1-u)}{u}$ gives $g(0^{+})=+\infty$ and $g(1^{-})=-\infty$. By the intermediate value theorem $g$ has a zero, unique since $g$ is strictly decreasing; call it $u_r\in(0,1)$. Thus $\psi$ increases on $[0,u_r]$ and decreases on $[u_r,1]$, and
\begin{equation}\label{eq:critical}
\kappa:=\psi(u_r)=h'(u_r).
\end{equation}
Set $\rho=e^{\kappa}$. From $h'(u)=\log\frac{r(1-u)}{u}$, equation \eqref{eq:critical} gives $r(1-u_r)/u_r=\rho$, that is $u_r=r/(r+\rho)$, and then $1-u_r=\rho/(r+\rho)$. Substituting into $h$,
\[
h(u_r)=u_r\log\tfrac{r+\rho}{r}+(1-u_r)\log\tfrac{r+\rho}{\rho}+u_r\log r
=\log(r+\rho)-(1-u_r)\log\rho ,
\]
while \eqref{eq:critical} reads $h(u_r)=\kappa(2+u_r)=(2+u_r)\log\rho$. Equating the two and cancelling,
\[
\log(r+\rho)=(2+u_r)\log\rho+(1-u_r)\log\rho=3\log\rho,
\]
that is
\begin{equation}\label{eq:cubic}
\rho^{3}=\rho+r .
\end{equation}
Since $z\mapsto z^{3}-z$ is strictly increasing for $z\ge1$ and vanishes at $z=1$, the root $\rho=\rho_r>1$ of \eqref{eq:cubic} is unique, and $\kappa=\log\rho_r$.

\emph{The two phases.} If $u_r\le u_{\max}(\beta)$, equivalently $\beta\le1/(1+u_r)=\beta^{*}$, the maximum is interior: $\Psi(\beta)=(1+\beta)\kappa=(1+\beta)\log\rho$, with unique maximizer $u=u_r$ and $\ell_\beta=(1+\beta)/(2+u_r)$. If $\beta>\beta^{*}$, the constraint binds, and since $\psi$ is strictly increasing on $[0,u_r]$ the unique maximizer is $u=(1-\beta)/\beta$; there $2+u=(1+\beta)/\beta$, so
\[
\Psi(\beta)=(1+\beta)\,\frac{h(u)}{(1+\beta)/\beta}=\beta\,h\Big(\tfrac{1-\beta}\beta\Big),
\qquad
\ell_\beta=\beta,\quad \mu_\beta=1-\beta .
\]
Finally $\beta^{*}=1/(1+u_r)=(r+\rho)/(2r+\rho)=\rho^{3}/(\rho^{3}+r)$ by \eqref{eq:cubic}. Uniqueness of the maximizer follows from the strict unimodality of $\psi$ together with \eqref{eq:partials}.
\end{proof}

\end{document}
